\documentclass[11pt]{article}
\usepackage{amsmath,amssymb,booktabs}
\title{Token Economics of Retention Policies}
\begin{document}
\maketitle
\section{Introduction}
Let $R$ denote the retention horizon and $D$ the disposition date. We write $\mathcal{P}$ for the set of policies and $\pi(d)$ for the policy resolved at document $d$.
\begin{equation}
  \pi(d) = \begin{cases} p_d & \text{if } p_d \neq \emptyset \\ \pi(\mathrm{parent}(d)) & \text{otherwise} \end{cases}
\end{equation}
\section{Cost model}
The expected storage cost is
\begin{align}
  C(\mathcal{P}) &= \sum_{d \in \mathcal{D}} \int_{t_0}^{t_0 + R(\pi(d))} c(t)\,\mathrm{d}t \\
                 &\leq |\mathcal{D}| \cdot \bar{c} \cdot \max_{p \in \mathcal{P}} R(p).
\end{align}
\begin{table}[h]
  \centering
  \begin{tabular}{lrr}
    \toprule
    Policy & Horizon (days) & Documents \\
    \midrule
    Default & 365 & 124000 \\
    Finance & 2555 & 18400 \\
    Legal & 3650 & 9100 \\
    Ephemeral & 30 & 412000 \\
    Audit & 1825 & 22300 \\
    \bottomrule
  \end{tabular}
  \caption{Retention horizons by policy.}
\end{table}
\section{Conclusion}
We have shown that $C(\mathcal{P})$ is minimised when $R$ is resolved at creation time rather than retroactively.
\end{document}
